%% slide13_study_note.tex
%% Detailed Study Note — Slide 13: Ch 9 Non-Blocking EKF Digital Twin
%% TSA Meeting Preparation | Anoop V Eluvathingal | IIT Madras | 2026
\documentclass[12pt,a4paper]{article}

\usepackage[margin=2.5cm]{geometry}
\usepackage{amsmath,amssymb}
\usepackage{booktabs,tabularx,array}
\usepackage{graphicx}
\usepackage{xcolor}
\usepackage{tcolorbox}
\usepackage{enumitem}
\usepackage{fancyhdr}
\usepackage{titlesec}
\usepackage{hyperref}
\usepackage{parskip}
\newcommand{\kibr}{k_\mathrm{ibr}}

\definecolor{iitmnavy}{RGB}{15,40,80}
\definecolor{iitmgold}{RGB}{180,140,0}
\definecolor{lightblue}{RGB}{220,235,250}
\definecolor{lightyellow}{RGB}{255,252,220}
\definecolor{lightred}{RGB}{255,235,235}
\definecolor{lightgreen}{RGB}{230,255,230}
\definecolor{lightpurple}{RGB}{240,230,255}

\tcbuselibrary{skins,breakable}
\newtcolorbox{keybox}[1]{colback=lightblue,colframe=iitmnavy,
  fonttitle=\bfseries\color{white},title=#1,arc=3pt,boxrule=1pt,breakable}
\newtcolorbox{formulabox}[1]{colback=lightyellow,colframe=iitmgold,
  fonttitle=\bfseries,title=#1,arc=3pt,boxrule=1pt,breakable}
\newtcolorbox{resultbox}[1]{colback=lightgreen,colframe=green!50!black,
  fonttitle=\bfseries,title=#1,arc=3pt,boxrule=1pt,breakable}
\newtcolorbox{warnbox}[1]{colback=lightred,colframe=red!60!black,
  fonttitle=\bfseries,title=#1,arc=3pt,boxrule=1pt,breakable}
\newtcolorbox{archbox}[1]{colback=lightpurple,colframe=iitmnavy!60,
  fonttitle=\bfseries,title=#1,arc=3pt,boxrule=1pt,breakable}
\newtcolorbox{speakbox}{colback=orange!8,colframe=orange!70!black,
  fonttitle=\bfseries,title={What to Say at the TSA Meeting},arc=3pt,
  boxrule=1pt,breakable}

\pagestyle{fancy}
\fancyhf{}
\fancyhead[L]{\color{iitmnavy}\small\textbf{TSA Meeting Study Note}}
\fancyhead[R]{\color{iitmnavy}\small Slide 13 — Non-Blocking EKF Digital Twin}
\fancyfoot[C]{\thepage}
\fancyfoot[R]{\small\color{gray}Anoop V Eluvathingal $\cdot$ IIT Madras $\cdot$ 2026}
\renewcommand{\headrulewidth}{1.5pt}
\renewcommand{\headrule}{\color{iitmnavy}\hrule width\headwidth height\headrulewidth}

\titleformat{\section}{\large\bfseries\color{iitmnavy}}{}{0em}{}[\color{iitmnavy}\titlerule]
\titleformat{\subsection}{\normalsize\bfseries\color{iitmnavy!80}}{}{0em}{}

\hypersetup{colorlinks=true,linkcolor=iitmnavy,urlcolor=iitmnavy}

\begin{document}

%% ── Title Block ──────────────────────────────────────────────────────────────
\begin{tcolorbox}[colback=iitmnavy,coltext=white,arc=4pt,boxrule=0pt]
\begin{center}
{\Large\bfseries TSA Meeting Preparation --- Detailed Study Note}\\[6pt]
{\large Slide 13: Ch 9 --- Non-Blocking EKF Digital Twin}\\[4pt]
{\normalsize SAMBP Framework $\cdot$ Anoop V Eluvathingal $\cdot$ IIT Madras $\cdot$ 2026}
\end{center}
\end{tcolorbox}

\vspace{0.5em}
\begin{center}
\small\textit{Sources: SAMBP TR-43, TR-44, TR-45 / 2026}
\end{center}

%% ── 1. Overview ──────────────────────────────────────────────────────────────
\section{1.\; Slide Overview and Purpose}

\begin{keybox}{What This Slide Is About}
This slide introduces the \textbf{Digital Twin (DT)} — the real-time brain of the
entire SAMBP framework. While every other slide describes a specific protection
element (87L, OC, distance, microgrid), this slide describes the \emph{shared
infrastructure} that feeds them all.

The DT continuously estimates three network parameters
$\boldsymbol{\theta} = \{k_\mathrm{ibr},\, Z_\mathrm{line},\, Y_\mathrm{shunt}\}$
from PMU and IED measurements. These estimates are broadcast to every protection
module in real time, enabling all adaptive thresholds and confidence gates to
respond to the current state of the IBR fleet — without human intervention.

The slide establishes three key properties of the DT: \textbf{accuracy} (RMSE
results), \textbf{speed} (convergence $<2$ cycles), and \textbf{safety}
(non-blocking and bounded update law).
\end{keybox}

\noindent\textbf{Position in the thesis:} Chapter~6 (Monitoring and ML), slide~13
of~22. This is the first of three Ch~6 slides (DT architecture, PINN classifier,
online learning). It is a pivotal slide because the DT underpins every contribution:
without reliable $\hat{k}_\mathrm{ibr}$, all adaptive thresholds revert to fixed settings.

%% ── 2. The Role ──────────────────────────────────────────────────────────────
\section{2.\; The Role of the Digital Twin}

\begin{archbox}{DT as the Central Parameter Hub}
\begin{center}
\begin{tabular}{@{}lp{9cm}@{}}
\toprule
\textbf{Module receives} & \textbf{What it uses from the DT}\\
\midrule
87L (Chapter~3) & $\hat{k}_\mathrm{ibr}$ → adaptive slope $k_m^*(\kibr)$\\
87T (Chapter~3) & $\hat{k}_\mathrm{ibr}$ → adaptive transformer bias $k_T^*(\kibr)$\\
SyncOC (Chapter~4) & $\hat{k}_\mathrm{ibr}$ → confidence gate threshold\\
Distance 21 (Chapter~4) & $\hat{Z}_\mathrm{line}$ → reach correction, IBR shift\\
Microgrid 67 (Chapter~5) & $\hat{k}_\mathrm{ibr}$ → mode detector prior\\
PINN (Chapter~6) & $\hat{k}_\mathrm{ibr}$ → feature input to classifier\\
WAPC (Chapter~7) & all $\boldsymbol{\theta}$ → wide-area settings update trigger\\
\bottomrule
\end{tabular}
\end{center}
Without the DT, every module above would need its own parameter estimator.
The DT centralises this into a single, validated, non-blocking component.
\end{archbox}

%% ── 3. EKF Prediction-Correction ────────────────────────────────────────────
\section{3.\; EKF Prediction--Correction Algorithm}

\subsection{The state vector}

The DT tracks three parameters jointly:
\begin{itemize}[nosep]
  \item $k_\mathrm{ibr}$ — IBR penetration ratio (primary, controls all thresholds)
  \item $Z_\mathrm{line}$ — line impedance magnitude (for distance reach and fault location)
  \item $Y_\mathrm{shunt}$ — shunt admittance (charging current, affects 87L stability)
\end{itemize}

\subsection{The update equation (from the slide)}

\begin{formulabox}{EKF Prediction--Correction (slide equation)}
\begin{equation}
  \hat{\boldsymbol{\theta}}_{k|k} = \hat{\boldsymbol{\theta}}_{k|k-1}
    + K_k\!\left(z_k - h\!\left(\hat{\boldsymbol{\theta}}_{k|k-1}\right)\right)
  \label{eq:ekf}
\end{equation}
\begin{itemize}[nosep]
  \item $\hat{\boldsymbol{\theta}}_{k|k}$ — updated state estimate at step $k$
        (posterior)
  \item $\hat{\boldsymbol{\theta}}_{k|k-1}$ — predicted state (prior, from
        previous cycle)
  \item $K_k$ — Kalman gain matrix: $K_k = P_{k|k-1}H_k^\top
        (H_k P_{k|k-1} H_k^\top + R_k)^{-1}$
  \item $z_k$ — measurement vector: PMU voltage phasors, IED current measurements,
        and the sequence current ratio $I_2/I_1$
  \item $h(\cdot)$ — nonlinear measurement model (hence \emph{Extended} KF — the
        Jacobian $H_k = \partial h/\partial\boldsymbol{\theta}$ is evaluated
        at each step)
  \item $z_k - h(\hat{\boldsymbol{\theta}}_{k|k-1})$ — innovation: difference between
        measured and model-predicted outputs; drives the update
\end{itemize}
\end{formulabox}

\subsection{Why EKF and not a simple RLS?}

Recursive Least Squares (used in TR-43) estimates $k_\mathrm{ibr}$ alone. The EKF
jointly estimates all three states, propagating their \emph{covariance} — it knows
how uncertain each estimate is and weighs the Kalman gain accordingly. This matters
because:
\begin{enumerate}[nosep]
  \item $Y_\mathrm{shunt}$ and $Z_\mathrm{line}$ are coupled in the measurement
        model; estimating them separately introduces bias.
  \item The EKF's covariance $P_{k|k}$ feeds directly into the confidence score
        $\gamma$ used by every adaptive protection module.
  \item Pre-fault warm-starting (tracking $\boldsymbol{\theta}$ in normal operation)
        reduces convergence time by \textbf{7$\times$} compared to a cold start.
\end{enumerate}

%% ── 4. Non-Blocking Property ─────────────────────────────────────────────────
\section{4.\; The Non-Blocking Property}

\begin{warnbox}{Why Non-Blocking Matters — The Safety Argument}
If the EKF update were \emph{in the critical path} of the relay trip decision, then:
\begin{itemize}[nosep]
  \item A slow convergence or numerical issue in the EKF could \textbf{delay a
        trip signal} by tens of milliseconds — causing a missed fault clearance.
  \item This would violate IEC~60255 operating time requirements and potentially
        cause equipment damage or arc flash.
\end{itemize}
The non-blocking design \emph{guarantees} this cannot happen.
\end{warnbox}

\subsection{Implementation: the shadow register pattern}

\begin{center}
\begin{tabular}{@{}lp{5cm}p{5.5cm}@{}}
\toprule
\textbf{Thread} & \textbf{Cycle time} & \textbf{Action}\\
\midrule
\textbf{Relay thread} (high priority) & 50~$\mu$s & Reads $\hat{\boldsymbol{\theta}}$
  from the shadow register; executes trip logic; issues GOOSE\\
\textbf{EKF thread} (low priority) & 20~ms (50~Hz) & Runs prediction--correction;
  writes result to shadow register when done\\
\midrule
\multicolumn{3}{l}{\textit{The relay thread never waits for the EKF thread.}}\\
\bottomrule
\end{tabular}
\end{center}

\subsection{The convergence guard (from the slide)}

The slide states: \emph{``if convergence $>3$ cycles, previous $\hat{\boldsymbol{\theta}}$
is held — relay never waits.''}

This means: if the EKF has not converged after 3 update cycles ($3 \times 20 = 60$~ms),
the relay continues using the last valid estimate. This is safe because:
\begin{itemize}[nosep]
  \item The pre-fault warm-start estimate is already within $\pm 5\%$ of the true value.
  \item The protection elements are designed with \emph{margin} that accommodates
        $\pm 10\%$ uncertainty in $\hat{k}_\mathrm{ibr}$.
  \item A 3-cycle convergence failure is a diagnostic flag, not a protection failure.
\end{itemize}

\subsection{End-to-end latency (TR-45)}

\begin{center}
\begin{tabular}{@{}lrr@{}}
\toprule
\textbf{Stage} & \textbf{Latency [ms]} & \textbf{Cumulative [ms]}\\
\midrule
CT/VT transducer          & 0.5  & 0.5 \\
IEC~61869-9 digital sampling & 0.5 & 1.0 \\
Merging unit $\to$ relay  & 1.0  & 2.0 \\
Relay algorithm (Z1/87L)  & 18.0 & 20.0 \\
GOOSE trip signal         & 2.0  & 22.0 \\
Circuit breaker operation & 30.0 & 52.0 \\
\midrule
\textbf{Total fault clearance} & \textbf{52.0} & \\
$T_\mathrm{mem}$ margin   & 48.0 & (within 100~ms) \\
\textbf{DT latency contribution} & \textbf{0.0} & \textbf{parallel, non-blocking} \\
\bottomrule
\end{tabular}
\end{center}

%% ── 5. Bounded Update Law ────────────────────────────────────────────────────
\section{5.\; Bounded Update Law}

\begin{formulabox}{Bounded Update Law (slide equation)}
\begin{equation}
  \|\boldsymbol{\theta}_{k+1} - \boldsymbol{\theta}_k\| \leq \delta_\mathrm{max}
  \label{eq:bounded}
\end{equation}
This constraint clips the EKF step size at each cycle. If the raw Kalman update
would produce a step larger than $\delta_\mathrm{max}$, the update is scaled back
to exactly $\delta_\mathrm{max}$ in the $\ell_2$ norm.
\end{formulabox}

\subsection{Why this is needed — the adversarial data threat}

A \textbf{CT-injection attack} (or a CT saturation transient) can cause a sudden
large spike in the measured currents, which would cause the EKF innovation
$z_k - h(\hat{\boldsymbol{\theta}}_{k|k-1})$ to become very large. Without the
bound, a single bad measurement could:
\begin{itemize}[nosep]
  \item Shift $\hat{k}_\mathrm{ibr}$ by 0.5~pu in one cycle, disabling the 87L element.
  \item Cause the distance relay to ``see'' a phantom fault location outside Zone~1.
\end{itemize}
The bounded update law limits the damage any single measurement can do to the
state estimate. This connects directly to Slide~17 (Ch~11 cybersecurity) where the
adversarial PINN robustness framework is formally established.

\subsection{Design of $\delta_\mathrm{max}$}

$\delta_\mathrm{max}$ is set to be:
\begin{itemize}[nosep]
  \item \emph{Large enough} to track genuine fast fleet changes
        (typical IBR dispatch ramp: $\Delta k_\mathrm{ibr} \approx 0.02$~pu/min
        $= 3.3 \times 10^{-4}$~pu/cycle at 50~Hz).
  \item \emph{Small enough} to reject adversarial step injections
        ($\delta_\mathrm{attack} \sim 0.1$--$0.5$~pu/cycle from CT saturation).
  \item In practice: $\delta_\mathrm{max} = 0.05$~pu provides a $150\times$ margin
        between tracking need and attack rejection.
\end{itemize}

%% ── 6. Results ───────────────────────────────────────────────────────────────
\section{6.\; Results — What the Table in the Slide Means}

\begin{resultbox}{EKF Digital Twin Results Table Explained}
\begin{tabular}{@{}lp{2cm}p{7.5cm}@{}}
\toprule
\textbf{Parameter} & \textbf{RMSE} & \textbf{Explanation}\\
\midrule
$\hat{k}_\mathrm{ibr}$ & \textbf{0.024~pu} &
  Primary result. On a typical $k_\mathrm{ibr} = 0.20$~pu fleet, this is
  12\% relative error — well within the $\pm 10\%$ protection margin.
  Maintained even at $>5\%$ PMU measurement noise.\\[4pt]
$\hat{Z}_\mathrm{line}$ (mag.) & \textbf{0.018~pu} &
  Line impedance magnitude. Lower RMSE than $k_\mathrm{ibr}$ because
  $Z_\mathrm{line}$ changes slowly (temperature drift only) and pre-fault
  warm-starting provides a very accurate prior.\\[4pt]
$\hat{Y}_\mathrm{shunt}$ & \textbf{0.031~pu} &
  Shunt admittance (charging current). Hardest to estimate — it appears
  only weakly in the current ratio measurement. Acceptable because 87L
  stability uses a $3\sigma$ threshold that accommodates 0.03~pu $Y_\mathrm{shunt}$
  uncertainty.\\[4pt]
Convergence time & \textbf{$<$2 cycles} &
  With warm-starting: EKF reaches $<3\%$ error within $1.9$ updates
  ($\approx$38~ms at 50~Hz). This is before Zone-1 distance relay operates
  (20~ms) only with the warm-start advantage; cold-start would take 13 updates.\\[4pt]
Max update latency & \textbf{1.8~ms} &
  Total DT cycle computation time within the relay firmware. This is the
  worst-case time the EKF thread holds the shadow register write lock —
  the relay thread is never blocked for longer than this.\\
\bottomrule
\end{tabular}
\end{resultbox}

\subsection{The headline observation}

\emph{``RMSE $\leq 0.024$~pu for $k_\mathrm{ibr}$ across all operating conditions
including $>5\%$ PMU measurement noise.''}

This is the \textbf{robustness claim}: even with severely degraded PMU measurements
(5\% noise is double the IEC class~1P CT error of 2.5\%), the DT maintains
sufficient accuracy for all downstream protection decisions. This directly
addresses a common examiner concern: \textit{``what happens when PMU data quality
is poor?''}

%% ── 7. Best Supporting Figure ────────────────────────────────────────────────
\section{7.\; Best Supporting Figure}

\begin{keybox}{Recommended Figure — TR-45 Figure 3 (Non-Blocking Timeline)}
\textbf{Source:} SAMBP TR-45/2026, Figure~3 (page~3)\\
\textbf{Title:} End-to-end fault clearance latency chain\\[4pt]
The figure is a horizontal timeline showing:
\begin{itemize}[nosep]
  \item \textbf{Relay chain} (top, orange): CT/VT (0.5~ms) $\to$ MU (1~ms) $\to$
        relay algorithm (18~ms) $\to$ GOOSE (2~ms) $\to$ CB (30~ms) $= 52$~ms total
  \item \textbf{DT thread} (bottom, green): 1~ms cycle running in parallel,
        \emph{contributing zero latency}
  \item $T_\mathrm{mem} = 100$~ms boundary and 48~ms margin clearly marked
\end{itemize}
\textbf{Why it is the best match:} The figure directly proves the slide's central
claim — ``Update runs in a parallel thread; relay never waits.'' It is
self-explanatory, has immediate visual impact, and pre-empts the most obvious
committee question (\textit{``does the digital twin slow down your relay?''}).
\end{keybox}

\noindent\textbf{Secondary figure (TR-43, Figure~1):} The 4-layer DT architecture
diagram (Physical $\to$ Relay $\to$ Digital Twin $\to$ Validation) provides context
for explaining the DT's role as a system-wide parameter hub.

\noindent\textbf{Tertiary figure (TR-44, Figure~3):} The DT cycle budget pie chart
(285~$\mu$s used, 71.5\% margin) explains the 1.8~ms max update latency result and
demonstrates computational headroom for future extensions.

%% ── 8. Connections Across the Thesis ────────────────────────────────────────
\section{8.\; Connections Across the Thesis}

\begin{center}
\begin{tabular}{@{}lp{10cm}@{}}
\toprule
\textbf{Chapter/Contribution} & \textbf{Connection}\\
\midrule
C1 — Differential (Ch~3) & $\hat{k}_\mathrm{ibr}$ drives slope $k_m^*$ and
  $k_T^*$; $\hat{Y}_\mathrm{shunt}$ sets charging current threshold $I_c$\\
C2 — OC/Distance (Ch~4) & $\hat{k}_\mathrm{ibr}$ gates confidence score $\gamma$;
  $\hat{Z}_\mathrm{line}$ corrects distance reach\\
C2 — Microgrid (Ch~5) & $\hat{k}_\mathrm{ibr}$ feeds the Bayesian mode detector
  prior (slide~12 equation)\\
C3 — PINN (Ch~6, next slides) & $\hat{k}_\mathrm{ibr}$ is the most important
  input feature (37~pp permutation importance drop)\\
C4 — WAPC (Ch~7) & CUSUM monitors the EKF's $\hat{k}_\mathrm{ibr}$ trajectory
  to detect fleet composition drift and trigger settings update\\
Appendix A & EKF Jacobian derivation and local observability proof\\
\bottomrule
\end{tabular}
\end{center}

%% ── 9. Equations Quick Reference ─────────────────────────────────────────────
\section{9.\; Equations Quick Reference}

\begin{formulabox}{Key Equations for This Slide}
\textbf{EKF prediction step:}
\begin{equation}
  \hat{\boldsymbol{\theta}}_{k|k-1} = \hat{\boldsymbol{\theta}}_{k-1|k-1},
  \quad P_{k|k-1} = P_{k-1|k-1} + Q
\end{equation}

\textbf{Kalman gain:}
\begin{equation}
  K_k = P_{k|k-1} H_k^\top \bigl(H_k P_{k|k-1} H_k^\top + R\bigr)^{-1}
\end{equation}

\textbf{EKF correction step (slide):}
\begin{equation}
  \hat{\boldsymbol{\theta}}_{k|k} = \hat{\boldsymbol{\theta}}_{k|k-1}
    + K_k\bigl(z_k - h(\hat{\boldsymbol{\theta}}_{k|k-1})\bigr)
\end{equation}

\textbf{Bounded update constraint (slide):}
\begin{equation}
  \|\boldsymbol{\theta}_{k+1} - \boldsymbol{\theta}_k\| \leq \delta_\mathrm{max}
  \quad (\delta_\mathrm{max} = 0.05\;\text{pu in deployment})
\end{equation}

\textbf{Cramér--Rao Lower Bound (accuracy floor):}
\begin{equation}
  \mathrm{RMSE}(\hat{\boldsymbol{\theta}}) \geq
    \sqrt{\mathrm{diag}\!\left[\mathbf{I}(\boldsymbol{\theta})^{-1}\right]}
\end{equation}
where $\mathbf{I}(\boldsymbol{\theta})$ is the Fisher information matrix.
The reported RMSE values (0.024, 0.018, 0.031~pu) are all within a factor of
$1.5\times$ of the CRLB, confirming near-optimal estimation efficiency.

\textbf{Warm-start convergence speedup (TR-44):}
\begin{equation}
  \frac{N_\mathrm{cold}}{N_\mathrm{warm}} = \frac{13.3\;\text{updates}}{1.9\;\text{updates}}
  \approx 7\times
\end{equation}
\end{formulabox}

%% ── 10. Speaker Script ───────────────────────────────────────────────────────
\section{10.\; TSA Meeting Speaker Script}

\begin{speakbox}
``Slide~13 introduces what I think of as the \textbf{nervous system} of the SAMBP
framework — the non-blocking EKF digital twin.

Every adaptive threshold I showed in the previous slides — the 87L slope, the OC
pickup, the Bayesian mode detector, the distance reach correction — they all depend
on knowing $k_\mathrm{ibr}$, the current IBR penetration level. The digital twin
estimates this continuously from PMU and IED measurements, and broadcasts it to all
modules.

The estimation uses an Extended Kalman Filter because the measurement model is
nonlinear. The filter jointly tracks three parameters: IBR penetration $k_\mathrm{ibr}$,
line impedance $Z_\mathrm{line}$, and shunt admittance $Y_\mathrm{shunt}$.

The most important design decision is that the EKF runs in a \textbf{parallel
thread}. The relay never waits for the EKF to finish. If the filter hasn't
converged after three cycles, the relay simply uses the last valid estimate. This is
the shadow register pattern — the relay reads a snapshot; the EKF writes to the
background copy without ever blocking the protection logic.

I also impose a \textbf{bounded update law}: no single EKF step can change the
estimate by more than $\delta_\mathrm{max}$. This prevents a CT injection attack
or a CT saturation transient from corrupting the state estimate in one cycle and
causing a relay mis-operation.

The results: RMSE of 0.024~pu for $k_\mathrm{ibr}$, 0.018~pu for $Z_\mathrm{line}$,
and 0.031~pu for $Y_\mathrm{shunt}$. These are maintained even with PMU noise above
5\% — double the rated class accuracy. Convergence takes less than two power
frequency cycles thanks to the pre-fault warm-start. And the maximum EKF thread
latency is 1.8~milliseconds, which contributes \textbf{zero additional delay} to
the 52-millisecond fault clearance chain.''
\end{speakbox}

%% ── 11. Anticipated Committee Questions ──────────────────────────────────────
\section{11.\; Anticipated TSA Committee Questions and Answers}

\begin{enumerate}
\item \textbf{Q: What happens if the EKF diverges during a fault?}\\
      A: Divergence is detected by monitoring the innovation sequence. If
      $\|z_k - h(\hat{\boldsymbol{\theta}}_{k|k-1})\|$ exceeds $3\sigma$ for
      two consecutive cycles, the EKF is re-initialised with the warm-start
      prior. The relay continues with the last valid estimate throughout.
      The bounded update law also limits how far a diverging estimate can drift.

\item \textbf{Q: The RMSE for $Z_\mathrm{line}$ is 0.018~pu — how does this
      affect distance protection accuracy?}\\
      A: The distance Zone~1 reach is set at 80\% of the line impedance with a
      $\pm$10\% margin. A 0.018~pu RMSE on a typical $Z_\mathrm{line} = 0.05$~pu
      line is 36\% relative error — this \emph{is} significant for impedance
      magnitude, but the product $\alpha Z_\mathrm{line}$ (fault location $\times$
      impedance) is better conditioned and the distance relay uses the corrected
      apparent impedance locus, not $Z_\mathrm{line}$ directly.

\item \textbf{Q: Why not use a particle filter or unscented KF instead of EKF?}\\
      A: The EKF Jacobian is analytically derived (Appendix A) and the nonlinearity
      is mild for $k_\mathrm{ibr} \in [0.03, 0.85]$~pu. The EKF update completes in
      17.6~$\mu$s; a particle filter with $N=100$ particles would require
      $\sim$1.76~ms — exceeding the 1.8~ms cycle budget. The UKF is a reasonable
      future extension but provides no significant accuracy gain for this problem.

\item \textbf{Q: How do you handle PMU data dropouts or communication latency?}\\
      A: The EKF prediction step runs even without a new measurement, using only the
      process model. If a PMU frame is missing, $\hat{\boldsymbol{\theta}}_{k|k}
      = \hat{\boldsymbol{\theta}}_{k|k-1}$ (no correction applied). This is standard
      KF missing-data handling and is safe for up to 200~ms (10 missed cycles at
      50~Hz) before the estimate uncertainty grows beyond the protection margin.

\item \textbf{Q: Is the bounded update law compatible with fast fleet changes,
      such as sudden IBR disconnection?}\\
      A: A sudden disconnection causes $k_\mathrm{ibr}$ to drop by up to 0.3~pu
      in one event. With $\delta_\mathrm{max} = 0.05$~pu/cycle, the EKF takes
      $0.30/0.05 = 6$ cycles ($\approx$120~ms) to track this change. During this
      window, the relay uses a slightly over-estimated $k_\mathrm{ibr}$, making
      protection elements more conservative (higher threshold). This is fail-safe:
      it produces fewer false trips rather than missed trips.
\end{enumerate}

%% ── Footer ───────────────────────────────────────────────────────────────────
\vspace{2em}
\begin{center}
\small\color{gray}
Study note prepared from SAMBP Technical Reports TR-43, TR-44, TR-45 (2026).\\
Best supporting figure: TR-45/2026, Figure~3 — End-to-End Non-Blocking Latency Timeline.\\
File: \texttt{slide13\_study\_note.pdf} $\cdot$ Generated \today
\end{center}

\end{document}
